%
%  Time-stamp: "[publicabstract.tex] last modified by Scott Budge (scott) on 2011-08-09 (Tuesday, 9 August 2011) at 09:17:43 on goga"
%
%  Info: $Id$   USU
%  Revision: $Rev$
% $LastChangedDate$
% $LastChangedBy$
%

\begin{publicabstract}
% A space is needed before the text starts so that the first paragraph
% is indented properly.

Snoring can wreck a good night's sleep, not just for the snorer but for their partner, and over time it can strain relationships. Common fixes like CPAP machines, nasal strips, and earplugs either have to be worn or don't do much against the deep, low-pitched rumble that makes snoring so hard to ignore.

This research explores a different approach, a device that doesn't need to be worn at all. It uses a small set of microphones to listen to the snoring and a set of speakers to play back carefully timed anti-snore noise, which is sound designed to cancel out the snore. The system focuses this cancellation on the area around the sleeping partner's ear, so the snoring is quieter where it matters while the rest of the room is left largely undisturbed.

In testing, the system reduced snoring noise at the target location, with the biggest improvements at low frequencies where traditional methods struggle. These results suggest that targeted, wearable-free noise cancellation could offer a practical way to help couples sleep better.

\end{publicabstract}


% Local Variables:
% TeX-master: "newhead"
% End:
